% sections/design.tex -- research-creation design, phases and participants (shared).
% Claims and sources: source_verification.md, sections A and E.
The project follows the program's research-creation option: a substantial
creative work and a 150--200-page written dissertation that carries the
theoretical argument \citep{comcult_rc}. Its five phases overlap and feed
back into one another (Figure~\ref{fig:overview}).

\textbf{Phase 1: Studio inquiry and tool ecology.} Before any participants
are involved, I will undertake a sustained period of my own artistic
research into collaboration and GenAI. I will select three or four tool
families, for example text, image, voice or audio, and spatial media
integrated through Unity, against explicit criteria: artistic relevance,
cost and access, data and privacy terms, whether inputs and outputs can be
documented, degree of user control, known cultural defaults and feasibility.
Where quality allows, one will be a locally run open-weight model, and
non-AI processes stay in the ecology. The phase produces experiments, scores,
a failure log, a reflexive journal and preliminary works, and reads
interface documentation, terms of service and model cards as cultural
documents. It asks how each tool structures prompting, refusal and
authorship; it is not software benchmarking.

\textbf{Phase 2: Consultation and pilots.} Four to six adult artists,
improvisers, facilitators or creative technologists will act as advisors in
pilot sessions that refine scores, access provisions, tool choices and
consent and credit protocols. They are not main participants, and I will
call them co-researchers only if they share decisions and credit.

\textbf{Phase 3: Collaborative workshops.} The provisional sample is 18--24
adults in three cohorts of six to eight: artists, performers and
improvisers, designers, creative technologists and community-engaged arts
practitioners. Open calls to arts communities in Toronto and York Region and
to AMPD networks, excluding anyone I teach, will seek variation in
discipline, GenAI experience, adult age and cultural and linguistic
background; the sample is not representative. I begin with adults because
the work is public, consent and authorship are complex to negotiate, and the
platforms' data practices are contested. Six to eight people can make work
together and be observed closely; three cohorts allow cross-case comparison. The scale rests on information power, meaning a
specific sample, an established framework and repeated, dense encounters
\citep{malterud2016}, not on statistical power or declared saturation
\citep{braun2021saturate}.

Each cohort will attend four connected workshops of about two and a half
hours. The sequence, set after Phases 1 and 2, will include: agreeing
group conventions and making together without assuming AI is needed; experiments with more than one form of GenAI mediation and with
AI/non-AI combinations; repeated chances to accept, revise, reroute, refuse,
remove or reinterpret outputs; collective rules for attribution, data use,
credit, reuse and decisions; reflection on misreading, standardization,
surprise, exclusion and productive failure; and movement toward materials
and scores for the final artwork. I will conduct disclosed participant
observation with field notes and the reflexive journal. With consent, I will
collect recordings, prompts and outputs, version histories, logs, artifacts
and group-authored rules. A log records system events, not intentions.

\textbf{Phase 4: Focus groups.} Each cohort will then join one
semi-structured focus group of 60--90 minutes on voice, recognition,
cultural legibility, the authority of models and platforms, ownership,
responsibility, refusal, access and desired changes. Focus groups make
participants' interaction part of the data \citep{kitzinger1994}. Short
individual follow-ups are offered only to those who need a private way to
discuss disagreement or power.

\textbf{Phase 5: Co-creation and exhibition.} The workshops lead toward a
large collaborative artwork proposed for the first-floor Art Gallery at
York's Markham Campus, which has shown creative-technologies exhibitions
\citep{yorkmarkham2024}, subject to supervisory, venue, resource, ethics and
participant approval. Its form (projection, spatial sound, interactive 3D
media, performance or installation) will emerge from the workshops. It
should make the research perceptible, not illustrate findings, for instance
by exposing traces of rejection, model defaults and contested credit. Participants decide whether and
how their material, names, voices and images appear; consent to exhibit is
separate from consent to research, and no audience data will be collected.
For examination, the work will be documented through installation records,
executable media where feasible, scores, process records and an annotated
account of its iterations.
